\documentclass[10pt,a4paper]{article}
\usepackage[utf8]{inputenc}
\usepackage[T1]{fontenc}
\usepackage[ngerman]{babel}
\usepackage[margin=1.7cm,top=1.5cm,bottom=1.7cm]{geometry}
\usepackage{titlesec}
\titleformat*{\section}{\normalfont\large\bfseries}
\titleformat*{\subsection}{\normalfont\normalsize\bfseries}
\titlespacing*{\section}{0pt}{8pt}{3pt}
\titlespacing*{\subsection}{0pt}{6pt}{2pt}
\usepackage{amsmath,amssymb}
\usepackage{listings}
\usepackage{xcolor}
\usepackage{fancyhdr}
\usepackage{parskip}
\usepackage{needspace}

\newcommand{\mcbox}{\raisebox{-1pt}{\fbox{\rule{0pt}{7pt}\rule{7pt}{0pt}}}\hspace{6pt}}
\newcommand{\antwortlinien}[1]{%
  \par\vspace{2pt}%
  \newcount\zeilen \zeilen=#1
  \loop
    \noindent\rule{\linewidth}{0.4pt}\par\vspace{11pt}%
    \advance\zeilen by -1
  \ifnum\zeilen>0 \repeat
}
\lstset{language=Python,basicstyle=\ttfamily\footnotesize,keywordstyle=\bfseries,
  commentstyle=\itshape\color{gray},showstringspaces=false,frame=single,
  framerule=0.3pt,xleftmargin=8pt,framexleftmargin=6pt,aboveskip=6pt,belowskip=6pt}

\pagestyle{fancy}
\fancyhf{}
\fancyhead[L]{\small MDT5/2 -- Session 10: Deep SVDD, GOAD, CutPaste}
\fancyhead[R]{\small Übungsaufgaben zur Klausurvorbereitung}
\fancyfoot[C]{\small Seite \thepage}
\renewcommand{\headrulewidth}{0.4pt}

\begin{document}

\begin{center}
  {\Large\bfseries Übungsaufgaben Session 10}\\[4pt]
  {\large Kapitel 8/8a: Deep SVDD, GOAD, CutPaste, selbstüberwachtes Lernen -- Praktika 10 + 11}\\[4pt]
  \small Stil der Übungssammlung MDT5/2 -- generierte Aufgaben, keine Originale
\end{center}

\vspace{4pt}

\section*{Aufgabe 1 -- Multiple Choice mit Begründung}
\emph{Markieren Sie jeweils die einzige richtige von drei Aussagen und begründen Sie Ihre Entscheidung.}

\subsection*{1.1 \ Für Deep SVDD gilt:}
\mcbox Das Kugelzentrum $\boldsymbol{c}$ wird zusammen mit den Gewichten per Gradient
Descent mitoptimiert.

\mcbox Das Kugelzentrum $\boldsymbol{c}$ wird vor dem Training einmalig als mittlerer
Vektor aller Abbildungen gesetzt und danach nicht mehr verändert.

\mcbox Das Kugelzentrum $\boldsymbol{c}$ wird auf den Nullvektor gesetzt, damit die
Optimierung einfacher wird.

Begründung:
\antwortlinien{2}

\subsection*{1.2 \ Für die Architektur einer Deep SVDD gilt:}
\mcbox Sigmoid-Aktivierungen in den Hidden Layers sind unproblematisch.

\mcbox Bias-Parameter und \glqq gedeckelte\grqq{} Aktivierungen müssen vermieden werden,
da sonst eine konstante (kollabierte) Abbildung möglich wird.

\mcbox Bias-Parameter sind zwingend erforderlich, damit das Netz flexibel genug ist.

Begründung:
\antwortlinien{2}

\subsection*{1.3 \ Für GOAD gilt:}
\mcbox GOAD trainiert ein Netz, die auf ein Sample angewendete Transformation zu erkennen;
pro Transformation soll ein dichtes Cluster im Latent Space entstehen.

\mcbox GOAD funktioniert ausschließlich auf Bilddaten.

\mcbox GOAD benötigt gelabelte Anomalien für das Training der Transformationen.

Begründung:
\antwortlinien{2}

\subsection*{1.4 \ Für selbstüberwachtes Lernen gilt:}
\mcbox Es benötigt manuell erstellte Labels für alle Trainingsdaten.

\mcbox Es erzeugt Pseudo-Labels über eine Hilfsklassifikationsaufgabe und führt damit
ein überwachtes Training durch, um gute Merkmale zu lernen.

\mcbox Es ist ein Spezialfall des Reinforcement Learning.

Begründung:
\antwortlinien{2}

\section*{Aufgabe 2 -- Deep SVDD: die Verbotsliste}

Nennen Sie \textbf{vier} Konstruktionsregeln für Deep SVDD aus der Vorlesung und begründen
Sie jede einzeln mit dem Kollaps-Argument (triviale Lösung / konstante Abbildung):

\antwortlinien{8}

\section*{Aufgabe 3 -- Deep-SVDD-Code beurteilen}

\emph{Die folgenden Code-Ausschnitte sollen als Deep-SVDD-Netze verwendet werden.
Beurteilen Sie jeden Ausschnitt auf mögliche Fehler in Hinblick auf Deep SVDD und
begründen Sie jeweils.}

\Needspace*{8cm}
\begin{lstlisting}
# Variante A
eps = 1e-6
wd = 1e-6
svdd_net = tf.keras.Sequential([
    tf.keras.layers.InputLayer((40,)),
    tf.keras.layers.Dense(units=25, activation='relu',
        kernel_regularizer=tf.keras.regularizers.l2(wd)),
    tf.keras.layers.Dense(units=15, use_bias=False, activation='relu',
        kernel_regularizer=tf.keras.regularizers.l2(wd)),
    tf.keras.layers.Dense(units=10, use_bias=False,
        kernel_regularizer=tf.keras.regularizers.l2(wd))
])
initialize_svdd_with_autoencoder(train_data, svdd_net)
transform_center = np.mean(svdd_net.predict(train_data), axis=0) + eps
\end{lstlisting}
\antwortlinien{3}

\Needspace*{8cm}
\begin{lstlisting}
# Variante B
eps = 1e-6
wd = 1e-6
svdd_net = tf.keras.Sequential([
    tf.keras.layers.InputLayer((40,)),
    tf.keras.layers.Dense(units=25, use_bias=False, activation=None,
        kernel_regularizer=tf.keras.regularizers.l2(wd)),
    tf.keras.layers.LeakyReLU(0.1),
    tf.keras.layers.Dense(units=15, use_bias=False, activation=None,
        kernel_regularizer=tf.keras.regularizers.l2(wd)),
    tf.keras.layers.LeakyReLU(0.1),
    tf.keras.layers.Dense(units=10, use_bias=False,
        kernel_regularizer=tf.keras.regularizers.l2(wd))
])
initialize_svdd_with_autoencoder(train_data, svdd_net)
transform_center = np.mean(svdd_net.predict(train_data), axis=0) + eps
\end{lstlisting}
\antwortlinien{3}

\Needspace*{8cm}
\begin{lstlisting}
# Variante C
eps = 1e-6
wd = 1e-6
svdd_net = tf.keras.Sequential([
    tf.keras.layers.InputLayer((40,)),
    tf.keras.layers.Dense(units=25, use_bias=False, activation='tanh',
        kernel_regularizer=tf.keras.regularizers.l2(wd)),
    tf.keras.layers.Dense(units=15, use_bias=False, activation='tanh',
        kernel_regularizer=tf.keras.regularizers.l2(wd)),
    tf.keras.layers.Dense(units=10, use_bias=False,
        kernel_regularizer=tf.keras.regularizers.l2(wd))
])
initialize_svdd_with_autoencoder(train_data, svdd_net)
transform_center = np.zeros((10,))
\end{lstlisting}
\antwortlinien{4}

\section*{Aufgabe 4 -- Deep SVDD: Ablauf und Score}

a) Das Training einer Deep SVDD findet in zwei Schritten statt. Welchen?
\antwortlinien{2}

b) Wie wird der Radius $R$ nach dem Training bestimmt, und welche Rolle spielt dabei
$\nu$? (Stichwort Quantil)
\antwortlinien{3}

c) Geben Sie die Formel für den Score eines unbekannten $\boldsymbol{x}$ an und
erklären Sie, wie daraus direkt eine Klassifikation entsteht (Vorzeichen).
\antwortlinien{3}

d) Warum minimiert die vereinfachte Zielfunktion
$\min_W \frac{1}{n}\sum_i \|\phi(\boldsymbol{x}_i;W)-\boldsymbol{c}\|^2 + \frac{\lambda}{2}\sum_l \|W^{[l]}\|^2$
implizit das Kugelvolumen, obwohl $R$ gar nicht mehr vorkommt?
\antwortlinien{3}

\section*{Aufgabe 5 -- GOAD}

a) Bei Bildern kombiniert GOAD Rotationen (4), Translationen (9) und horizontale
Spiegelung (2). Wie viele Transformationen $M$ ergeben sich, und warum ist die Identität
enthalten?
\antwortlinien{2}

b) Der Triplet Center Loss besteht aus zwei gegenläufigen Zielen. Welchen?
(Intra- vs. Inter-Cluster)
\antwortlinien{3}

c) Warum wird zusätzlich eine Klassifikationsaufgabe (Kreuzentropie) als Stabilisierung
verwendet, und was wird dabei klassifiziert?
\antwortlinien{3}

d) Beschreiben Sie, wie der Anomalie-Score eines neuen Samples entsteht
(Wahrscheinlichkeiten pro Transformation, negativer Logarithmus).
\antwortlinien{4}

e) Wie wird GOAD auf Nicht-Bilddaten verallgemeinert?
\antwortlinien{2}

\section*{Aufgabe 6 -- CutPaste und Contrastive Learning}

a) Welches Problem der bisherigen Verfahren motiviert CutPaste?
\antwortlinien{2}

b) Beschreiben Sie die drei Klassen der CutPaste-Hilfsklassifikationsaufgabe.
\antwortlinien{3}

c) Wie entsteht bei CutPaste nach dem Training der Anomalie-Score
(welches Verfahren aus einem früheren Kapitel steckt dahinter)?
\antwortlinien{3}

d) Was ist die Grundidee des Contrastive Learning (positive/negative Samples,
Ähnlichkeitsmaß), und wie erzeugt SimCLR seine positiven Paare?
\antwortlinien{4}

\vspace{10pt}
\begin{center}
\footnotesize\itshape
Nach Bearbeitung: Antworten an Claude schicken (Foto genügt) -- Korrektur mit Begründungs-Check.
\end{center}

\end{document}
